\section{The Cut Happens First}
\label{sec:ch13-cut-happens-first}
\index{query plan!page cut before hydration}%

The plan explains the result without any special case for pagination. The
router first asks Sessions for one row. It can then build a Session
representation and ask Ratings for the score of that row. By the time Ratings
answers, membership has already been decided.

The two paths in figure~\ref{fig:ch13-two-page-orders} separate operations that
a flat GraphQL response makes easy to confuse. Hydration fills fields on chosen
objects. Discovery chooses the objects. The router is good at the first. It
does not invent the second.

\begin{figure}[htbp]
  \centering
  % Two page orders. In both lanes the page has two sessions and crosses the
% same subgraph seam twice: keys travel out and fields travel back. The five
% ticks are aligned because the operations have the same broad shape. What
% changes is which service owns enough data to order the list before it cuts
% the page.
%
% Lane A starts from the Sessions list. Sessions orders that list and applies
% first: 2 before the router asks Ratings for scores. Ratings can enrich only
% the two representations it receives. Once those scores return, the router
% can only merge fields into the two fetched rows. It cannot replace either row
% with a higher-scored session whose key was never in the request.
%
% Lane B starts from a Search index that contains the fields needed for both
% filtering and ordering. Search can therefore choose the two-session page
% while every value that determines membership is present. Sessions receives
% only the chosen keys and supplies canonical fields; the router's last merge
% does not have to revisit which sessions belong to the page.
%
% Same idiom as figures/tikz/ch01-one-field.tex and
% ch05-the-join-moves.tex: a left-to-right timeline in each lane, square
% black-on-white engineering labels, one stroke weight, and dashed vertical
% rules for the fixed service seam. The seam is crossed once in each direction,
% so each lane has two rules. This is deliberately not a box-and-arrow
% architecture diagram: sequence, and specifically the position of the cut,
% is the claim.
%
% Bare tikzpicture (house style): the figure environment, caption and label
% live at the call site in the section file.
\begin{tikzpicture}[
  x=1mm, y=1mm,
  every node/.append style={font=\sffamily\scriptsize},
  axis/.style={draw, line width=0.4pt,
    -{Straight Barb[length=1.4mm, width=1.6mm]}},
  tick/.style={draw, line width=0.4pt},
  event/.style={anchor=south, align=center, text width=24mm,
    inner sep=1pt},
  span/.style={draw, line width=0.4pt},
  spanlabel/.style={anchor=north, align=center, inner sep=2pt},
  lane/.style={anchor=west, font=\sffamily\scriptsize\itshape},
  boundary/.style={draw, line width=0.4pt, dashed},
]

% ------------------------------------------------------------------- lane A
\node[lane] at (0,21)
  {A. Sessions owns the list: scores arrive after the cut};

\draw[axis] (0,0) -- (150,0);
\foreach \x in {14,45,76,108,140}{
  \draw[tick] (\x,-1.5) -- (\x,1.5);
}
\foreach \x in {94,123}{
  \draw[boundary] (\x,-2) -- (\x,10);
  \node[spanlabel, inner sep=1pt] at (\x,-2) {seam};
}

\node[event, text width=22mm] at (14,3)
  {Sessions owns\\the list};
\node[event, text width=28mm] at (45,3)
  {Sessions orders,\\then cuts\\{\ttfamily first: 2}};
\node[event, text width=27mm] at (76,3)
  {router sends only\\those 2 session\\keys to Ratings};
\node[event, text width=23mm] at (108,3)
  {Ratings adds\\scores to those\\2 sessions};
\node[event, text width=25mm] at (140,3)
  {router cannot admit\\a higher-scored\\session it never\\fetched};

\draw[span] (35,-8) -- (35,-11) -- (55,-11) -- (55,-8);
\node[spanlabel] at (45,-12) {page membership fixed here};

% ------------------------------------------------------------------- lane B
\begin{scope}[shift={(0,-55)}]
  \node[lane] at (0,21)
    {B. Search owns a combined index: rank is known before the cut};

  \draw[axis] (0,0) -- (150,0);
  \foreach \x in {14,45,76,108,140}{
    \draw[tick] (\x,-1.5) -- (\x,1.5);
  }
  \foreach \x in {94,123}{
    \draw[boundary] (\x,-2) -- (\x,10);
    \node[spanlabel, inner sep=1pt] at (\x,-2) {seam};
  }

  \node[event, text width=22mm] at (14,3)
    {Search owns the\\combined index};
  \node[event, text width=28mm] at (45,3)
    {Search filters,\\orders, and cuts\\{\ttfamily first: 2}};
  \node[event, text width=27mm] at (76,3)
    {router sends the\\2 chosen session\\keys to Sessions};
  \node[event, text width=23mm] at (108,3)
    {Sessions supplies\\canonical fields\\for those 2};
  \node[event, text width=25mm] at (140,3)
    {router merges the\\chosen page; all\\required values were\\present at the cut};

  \draw[span] (35,-8) -- (35,-11) -- (55,-11) -- (55,-8);
  \node[spanlabel] at (45,-12) {page membership fixed here};
\end{scope}

\end{tikzpicture}

  \caption{In the measured plan, hydration cannot repair a page cut against
  the wrong ordering. The list owner must see the ordering values first.}
  \label{fig:ch13-two-page-orders}
\end{figure}

That distinction is visible in the measured plans for this graph. A root
\code{Single} fetch produces a list. A later \code{BatchEntity}
fetch receives representations for members of that list and fills a path such
as \code{sessions.nodes.@.averageScore}. There is no plan node that discards
members, reorders them, creates new cursors, and asks the root owner for a
replacement row.

The missing replacement matters. Imagine Sessions returns ten rows and the
router removes six whose scores are too low. Returning four rows violates
\code{first: 10}. Fetching ten more is not generally sufficient because the
next ten might also fail. The router would need a scan protocol, stable sort
semantics, cursor ownership, and a termination rule spanning two services.
None of that follows from selecting a field.

For a connection with this plan shape, I use a strict rule: its owner controls
membership, filtering, ordering, cursors, and the slice. Remote entity fields
may decorate the members. They do not participate in choosing them.
